\documentclass[12pt,a4paper]{article}
\usepackage[T1]{fontenc}
\usepackage[utf8]{inputenc}
\usepackage[magyar]{babel}
\usepackage{geometry}
\geometry{margin=2.5cm}

\title{Tartalomjegyzék vázlat}
\author{}
\date{}

\begin{document}

\maketitle
\tableofcontents
\newpage

\section{Rövidítés- és jelölésjegyzék}

\section{Bevezetés}
\subsection{Motiváció és a témaválasztás indoklása}
\subsection{A megoldandó feladat rövid, közérthető leírása}
\subsection{A dolgozat célja és hozzájárulásai}
\subsection{A dolgozat felépítése}

\section{Felhasználói dokumentáció}
\subsection{A megoldott probléma rövid megfogalmazása}
\subsection{A felhasznált módszerek rövid áttekintése}
\subsection{Rendszerkövetelmények}
\subsection{Telepítés és első indítás}
\subsection{A grafikus felhasználói felület bemutatása}
\subsection{Tipikus felhasználási forgatókönyvek}
\subsection{Hibaüzenetek és gyakori problémák}

\section{Fejlesztői dokumentáció}
\subsection{A probléma részletes specifikációja}
\subsection{Tervezési döntések és indoklásuk}
\subsection{Elméleti háttér és a használt fogalmak definíciója}
\subsection{Ellenfél alapú támadások áttekintése}
\subsection{FGSM (Fast Gradient Sign Method)}
\subsection{PGD (Projected Gradient Descent)}
\subsection{Kapcsolódó munkák és további támadási módszerek}
\subsection{Védekezési stratégiák áttekintése}
\subsection{A rendszer architektúrája}
\subsection{A felhasznált technológiák}
\subsection{Az egyes komponensek megvalósítása}

\section{Tesztelés}
\subsection{Tesztelési terv}
\subsection{Egységtesztek}
\subsection{Integrációs tesztek}
\subsection{Funkcionális tesztek és felhasználói validáció}
\subsection{Teszteredmények és értelmezésük}
\subsection{A tesztelés korlátai}

\section{Összefoglalás}
\subsection{Elért eredmények}
\subsection{Következtetések}
\subsection{Továbbfejlesztési lehetőségek}

\section{Mellékletek}

\end{document}
